\begin{textsample}{POS Dim 6 – vem\_ed\_32 – Score 53.00 – vem\_ed\_32/t167.txt}  \label{ex:f6_pos_005}
Jornada dos PCIs 2025: Abertura No dia 9 de outubro ocorreu a Jornada dos Projetos Colaborativos Internacionais 2025.
O webinário online e gratuito é realizado anualmente desde 2021, pela Coordenadoria Geral do Ensino Superior de Graduação (CGESG) do CPS.
A edição deste ano abordou o tema “Internacionalização, pesquisa e extensão”.
O objetivo da Jornada dos Projetos Colaborativos Internacionais é debater questões relativas aos Intercâmbios Virtuais, com foco nos Projetos Colaborativos Internacionais (PCIs) desenvolvidos nas Fatecs e em projetos do tipo COIL (Collaborative Online International Learning – Aprendizagem Colaborativa International realizada on-line) ou BRaVE (Brazilian Virtual Exchange), desenvolvidos em instituições que utilizam abordagens similares, independentemente da nomenclatura adotada.
A \textbf{programação} contou com \textbf{palestrantes} da Universidade Federal de Santa Catarina, Universidade Federal Fluminense e Universidad Nacional de Río Cuarto (Argentina).
E \textbf{apresentações} de Projetos Colaborativos Internacionais (PCIs) realizados nas Fatecs de Pompeia, Sorocaba, Santo André e Bebedouro.
A quinta edição da Jornada dos PCIs registrou 142 inscrições.
Na pesquisa realizada após o evento com os participantes, 86 \% dos \textbf{respondentes} afirmaram estar “muito \textbf{satisfeitos}” com a organização.
As \textbf{gravações} estão disponíveis na playlist do \textbf{canal} da CGESG no YouTube: https://bit.ly/4rpWGPD O número 27 de VEm resume a Jornada dos Projetos Colaborativos Internacionais 2024 e traz um histórico das edições anteriores do evento: Jornada dos Projetos Colaborativos Internacionais 2024 - https://tinyurl.com/3a7hcpr3 Após a abertura da Jornada dos PCIs 2025, realizada por Otávio de Moraes, chefe de \textbf{gabinete} da presidência do Centro Paula Souza, houve uma breve \textbf{apresentação} sobre a equipe de Apoio à Internacionalização da Divisão de Extensão e Pesquisa no Ensino Superior da CGESG e sobre a estruturação dos PCIs, seu histórico e crescimento.
Patrícia Patrício, responsável pela comunicação da equipe, e Osvaldo Succi Junior, coordenador dos PCIs, comentaram os números e a \textbf{evolução} dos projetos, envolvendo pesquisa e extensão.
Entre 2013 e 2025, foram realizados 750 PCIs envolvendo 57 Fatecs e 94 Instituições de Ensino Superior internacionais, engajando cerca de 15 mil fatecanos e 15 mil estudantes das instituições parceiras.
Por ano, são realizados cerca de 120 PCIs nas Fatecs, o que \textbf{coloca} o Centro Paula Souza entre as cinco maiores instituições do mundo em número de projetos de Intercâmbio Virtual do tipo COIL.
Palestras Luciane Stallivieri (Universidade Federal de Santa Catarina) proferiu a pa- lestra \textbf{magistral} “Pontes entre Internacionalização, Extensão e Pesquisa”.
Defendeu a internacionalização como eixo transversal que \textbf{articula} ensino, pesquisa e extensão. “Não são polos isolados, e sim um sistema integrado onde intercâmbios virtuais atuam como conectores práticos entre ensino, comuni- de e \textbf{investigação}”.
A professora apresentou \textbf{definições} de Internacionalização em Casa, Internacionalização do Currículo e Internacionalização da Educação Superior, Intercâmbio Virtual e COIL, entre outros conceitos relacionados.
E deu \textbf{sugestões} para que os Intercâmbios Virtuais ajudem a atingir objetivos extensionistas, gerem pesquisas e \textbf{promovam} internacionalização.
Cíntia Rabello (Universidade Federal Fluminense) \textbf{relatou} o “International Vir- tual Exchange Project e suas contribuições para a aprendizagem de língua in- glesa e letramentos digitais de licenciandos na Universidade Federal Fluminense”.
Trata-se de uma iniciativa financiada pelo governo do Japão desde 2015.
A professora participou com seus alunos de inglês na UFF, em cinco edições do projeto, entre 2019 e 2025. “É uma oportunidade de expandir a aprendizagem de língua inglesa para além da sala de aula por meio da comunicação \textbf{autêntica} entre estudantes de diferentes culturas”, ressalta Cíntia. “Apesar do engajamento e interesse, alguns alunos têm dificuldade de realizar as interações fora da sala de aula”, observa.
María Soledad Fontana e Silvia Cristina Bertolo (Universidad Nacional de Río Cuarto, Argentina) compartilharam uma experiência COIL trilíngue, entre 16 alunos de \textbf{Português} e Francês, \textbf{orientados} por quatro docentes da universidade argentina e 18 estudantes de Relações Internacionais da Unesp, \textbf{supervisionados} por dois professores.
O projeto “Restaurante internacional: jantares temáticos e eventos culturais no Brasil, Argentina e França” consistia na \textbf{elaboração} de um flyer em \textbf{português}, espanhol e francês com o menu de um evento \textbf{multicultural} relacionando os três países.
As produções foram publicadas no \textbf{Padlet}.
Apresentações Vânia Regina Alves de Souza, da Fatec Pompeia, discorreu sobre os PCIs realizados na \textbf{unidade} com Symbiosis Centre for Management Studies (SCMS / Índia), Yerevan State University (Armênia), Eastern Oregon University (EOU / EUA) e Bukidnon State University (Filipinas).
O projeto com a Índia envolve também a Fatec Americana, aborda temas como satisfação e motivação no trabalho e já está na terceira edição.
O PCI com Armênia ocorre desde 2023, engajando também as Fatecs Araçatuba, Bragança Paulista, Carapicuíba, Guaratinguetá, Jales, Sebrae e Tatuí.
Com Eastern Oregon University, o projeto na área de estudos comparativos sobre agricultura no Brasil e EUA ocorreu no primeiro semestre de 2025, com participação das Fatecs Pompeia e Mogi das Cruzes.
Os alunos dos EUA estudaram os biomas brasileiros e os alunos brasileiros pesquisaram as regiões norte-americanas.
O PCI com a universidade das Filipinas começou em outubro de 2025.
Como principais \textbf{benefícios} dessas experiências, Vânia citou a \textbf{melhoria} na interação com os alunos, a participação em eventos internacionais, escrita de artigos acadêmicos (por professores e alunos) e desenvolvimento de competências socioemocionais (\textbf{soft} skills).
Wanderley do Prado \textbf{relatou} sua primeira experiência com PCIs, que ocorreu entre Fatec São Roque e DUOC UC (Chile) no segundo semestre de 2024, sobre \textbf{elaboração} de anúncios utilizando memes.
No primeiro semestre de 2025, o professor iniciou um projeto, realizado em inglês, entre Fatec Sorocaba e UNAM, agora na segunda edição.
Cada grupo propõe um produto ou serviço a ser oferecido no Brasil e no México, abordando Objetivos de Desenvolvimento Sustentável (ODS).
Entre as principais motivações dos alunos, Wanderley identifica a prática do idioma estrangeiro em um ambiente interativo e divertido, além de conhecer novas pessoas e culturas.
Para os professores, perceber que o interesse dos alunos pelas aulas \textbf{aumentou}, publicar artigos com os parceiros internacionais sobre a experiência dos PCIs e participar de projetos de extensão institucional por meio dos PCIs.
Quanto ao idioma, o professor observou a diferença entre os alunos da Fatec São Roque, que têm aulas de inglês no currículo, e os de Sorocaba, que não têm (e \textbf{acabaram} buscando cursos fora da faculdade). “Quando o professor Osvaldo Succi me procurou \textbf{perguntando} sobre a possibi- lidade de um Projeto Colaborativo Internacional com a DUOC do Chile, eu tive a ideia de trocar informações: os \textbf{chilenos} nos apresentam os \textbf{métodos} de \textbf{homologação} \textbf{veicular} e eu trabalho com eles a sistemática de medições de potência, \textbf{torque} e \textbf{emissões} de \textbf{poluentes}.
O professor Orlando Salvo Junior, da Fatec Santo André, trabalha com essa área de inspeção \textbf{veicular} e \textbf{homologação}, que é um tema novo para nós.
Ler \textbf{norma} é uma coisa, ver acontecendo na prática é completamente diferente”.
Outra vantagem apresentada pelo coordenador foi incluir o PCI como projeto de extensão institucional. “Isso foi realizador”.
Além disso, Fróes destacou o engajamento dos alunos, que se ajudaram na pro- dução dos vídeos dos \textbf{ensaios} técnicos com os \textbf{veículos}.
A atividade reduziu reprovação por falta e \textbf{nota} na disciplina.
A \textbf{nota} média final subiu de 6,6 para 7,4 e a assiduidade \textbf{melhorou} 17 \%.
Selma de Fátima Grossi, coordenadora do curso de Logística, e Lígia de Gran- di, professora de espanhol na Fatec Bebedouro, falaram sobre a sinergia entre docentes do idioma espanhol e da \textbf{logística} em um PCI realizado com a Unimi- nuto (Colômbia).
O projeto está em sua quarta edição e aborda o tema “Logística \textbf{empresarial} e sustentável”, estudando aspectos relacionados à produção de \textbf{café}, banana, flores, carne e \textbf{cana-de-açúcar}, importantes na economia dos dois países.
A pri- meira edição (início de 2024) \textbf{enfocou} a temática da \textbf{logística} verde na cadeia de \textbf{suprimentos}.
No segundo semestre de 2024, a questão foi “como reduzir o impacto ambiental” desses produtos representativos.
No início de 2025, o tema foi “impacto ambiental do transporte terrestre” (rodoviário, que prevalece nos dois países).
No segundo semestre de 2025, a questão é “Do lixo ao \textbf{consumo}” (economia circular).
As professoras ressaltaram que o PCI permitiu conhecer aspectos do país \textbf{vizinho}, \textbf{melhorar} no idioma e desenvolver empatia.

% matched lemmas: acabar, apresentação, articular, aumentar, autêntico, benefício, café, cana-de-açúcar, canal, chileno, colocar, consumo, definição, elaboração, emissão, empresarial, enfocar, ensaio, evolução, gabinete, gravação, homologação, investigação, logística, magistral, melhorar, melhoria, multicultural, método, norma, nota, orientados, padlet, palestrante, perguntar, poluente, português, programação, promor, relatar, respondente, satisfeito, soft, sugestão, supervisionar, suprimento, torque, unidade, veicular, veículo, vizinho
\end{textsample}
